\section{Introduction}
\label{sec:introduction}

Static benchmarks for evaluating large language models (LLMs) become contaminated through ordinary use~\cite{sainz2023contamination}. Once a benchmark is published, its contents leak into future training corpora. This happens directly, when the published set is dragged from the web, and indirectly, when problems sent to commercial model APIs during evaluation are kept and reused for training. When a model sees a benchmark's problems, or modified versions of them, during training, its score displays recall, not the model's capability the benchmark was built to test. Sainz~\emph{et al.}~\cite{sainz2023contamination} argue that contamination should be measured for every benchmark; Balloccu~\emph{et al.}~\cite{balloccu2024leakcheat} quantify the leak, finding that roughly 42\% of 255 reviewed evaluation papers leaked data to closed-source models, some 4.7 million samples across 263 benchmarks. % cite:external -- "42" and "4.7" are Balloccu et al.'s own reported figures, verified directly against the source paper, not this project's computed values

Test generation for software is a specific instance of this problem. An LLM asked to write unit tests for a function could seem highly capable if it memorised that function's tests during training. The same model may fail on a function it has never seen. UnLeakedTestBench (ULT)~\cite{huang2026ult} is the closest most recent static baseline. Its authors decontaminate a set of real-world functions and pair each with a deliberately leaked twin, PLT, so any contamination signal becomes visible. But a decontaminated benchmark, once published, faces the same risk it was built to avoid however carefully it was decontaminated. A benchmark that generates fresh problem instances at evaluation time, from a private template grammar, avoids this.

This work builds such a benchmark and assesses that claim. DyMuT is built from sixty parameterised templates spanning four cyclomatic-complexity (CC) tiers. Each template produces distinct, seeded problem instances, scored against a reference oracle by a mutation-testing suite instead of a fixed set of expected outputs. The aim is to establish whether this design resists the score inflation that contamination causes in static benchmarks, while still separating strong test generation from weak.

Three objectives follow from this aim. First: Compare mutation-testing performance across nine models on DyMuT, ULT, and PLT, and check if either static benchmark shows the contamination pattern (higher scores on the leaked variant). Second: Fine-tune a small open-source model on ULT's ground truth, then measure whether the gain carries over to DyMuT or stays on the static benchmark to distinguish the real capability from inflation caused by contamination. Third: Construct validity by checking whether DyMuT's ranking of models agree with an independently maintained measure of code-generation ability. Fourth: Establish DyMuT's own measurement properties -- whether mutation-killing effectiveness discriminates cleanly by complexity tier, and whether its model ranking is stable across independent re-instantiations.

Five research questions directly address the objectives. RQ1 asks whether DyMuT resists the contamination a static benchmark absorbs; RQ2 asks how mutation-killing effectiveness changes with complexity tier; RQ3 asks whether any static-dynamic score gap is larger at low tiers than high ones; RQ4 asks whether DyMuT's model rankings are stable across independent re-instantiations; RQ5 asks whether DyMuT's ranking agrees with an external, independently maintained code-generation proxy. Six hypotheses test these questions.

This work's contributions are: DyMuT itself (Section~\ref{sec:methods}); a causal demonstration that contamination inflates a static benchmark without correspondingly inflating DyMuT's (Section~\ref{sec:results-h2}); a construct-validity check against an independent proxy (Section~\ref{sec:results-h6clause1}); and a disclosed inconsistency, noticed during fine-tuning preparation, between ULT's published codebase and its paper's methodology (Section~\ref{sec:discussion-ult}).

The scope is bounded: every problem is a single Python function; mutation testing uses seven semantic-altering-change operators only, semantic-preserving-change operators being out of scope (Section~\ref{sec:conclusion}); every result uses one sampled attempt per problem ($K=1$); the causal fine-tuning result uses the smallest open-source model, and one fine-tuning method, LoRA, so it is not claimed to generalise to a larger model or full fine-tuning. Section~\ref{sec:related-work} reviews prior work; Section~\ref{sec:methods} describes DyMuT's construction; Section~\ref{sec:results} reports the experimental runs; Section~\ref{sec:discussion} discusses the findings; Section~\ref{sec:ethics} considers ethical dimensions; Section~\ref{sec:conclusion} concludes.

\input{table1}
